\chapter{Building from the source}\label{app:building}

\RITMO{} is built with CMake and Qt. The full instructions are in the file \texttt{BUILD.md} of the source; here is the short version.

\section{Linux}

\begin{lstlisting}
sudo apt install cmake qt6-base-dev portaudio19-dev librtmidi-dev
cmake -B build-qt -DCMAKE_BUILD_TYPE=Release
cmake --build build-qt -j
./build-qt/out/ritmo song.rmt
\end{lstlisting}

The program is \texttt{build-qt/out/ritmo}, with the \texttt{resources} folder next to it. \texttt{-DRMT\_QT\_MAJOR=5} builds with
Qt~5.15 instead of Qt~6.

\section{Windows and macOS}

Windows is built with MSYS2 (MinGW64: GCC, Qt~6, PortAudio, RtMidi) and macOS with the Apple compiler and the Qt of the official installer. The
GitHub workflows of the repository (\texttt{.github/workflows}) build the packages of both, and the AppImage for Linux
(\texttt{scripts/build-appimage.sh}), and are the best description of the steps.

\section{Linux riscv64 and ppc64}

There is no AppImage for these architectures (no official Qt~6, no \texttt{linuxdeploy}). The workflow
\texttt{.github/workflows/build-linux-ports.yml} builds the program in a Debian container of the architecture emulated by QEMU:
\texttt{debian:trixie} for riscv64, a \texttt{debootstrap} of Debian ports \emph{sid} for ppc64. Two scripts do the work, and can be
started by hand in any Debian of the architecture:
\begin{lstlisting}
docker run --rm --platform linux/riscv64 -v "$PWD":/src -w /src debian:trixie \
    sh -c 'scripts/build-tarball.sh && scripts/build-deb.sh'
\end{lstlisting}
\texttt{scripts/build-tarball.sh} makes the \texttt{.tar.gz} with the libraries (\texttt{lib}, \texttt{plugins}) and the launcher;
\texttt{scripts/build-deb.sh} makes the \texttt{.deb}, whose dependencies come from \texttt{dpkg-shlibdeps}. Both run the program
without a display as a test.

The engine on PowerPC 64 big endian is tested on every push (\texttt{.github/workflows/test-bigendian.yml}):
\texttt{RmtCoreTest} is built for x86\_64 and cross-compiled for ppc64, run by \texttt{qemu-ppc64}, and
\texttt{scripts/test-endian.sh} compares the POKEY registers, the sound and the screen of every song of \texttt{rmt/songs}.

\section{Core only and tests}

\texttt{-DRMT\_BUILD\_CORE\_ONLY=ON} builds the engine and the audio and MIDI backends without Qt, and with \texttt{-DRMT\_CORE\_TEST=ON}
the \texttt{RmtCoreTest} program, a smoke test of the engine. The \texttt{scripts/test-scripting.sh} script runs the reference scripts.

\section{Test hooks}

For tests and documentation the program can run without a display (\texttt{QT\_QPA\_PLATFORM=offscreen}). The environment variables
\texttt{RMT\_QT\_GRAB} (save the window and quit), \texttt{RMT\_QT\_KEYS}, \texttt{RMT\_QT\_COMMANDS}, \texttt{RMT\_QT\_FILEDIALOG},
\texttt{RMT\_QT\_DIALOG} and \texttt{RMT\_QT\_MENU\_GRAB} (the menus as images) press keys, run menu commands, answer the dialogs and save them. The
screenshots of this manual were taken that way, by \texttt{doc/manual/make-screenshots.sh}.
